% ---------------------------------------------------------------------
% Front matter. Two versions controlled by \ifanonymous in main.tex.



\begin{titlepage}
\centering
\vspace*{2cm}
{\LARGE\bfseries \thetitle \par}
\vspace{2cm}

\ifanonymous
  {\large Exam number: \examnumber \par}
  \vspace{0.5em}
  {\large MSc Economics Dissertation \par}
  \vspace{0.5em}
  {\large \today \par}
  \vspace{2cm}
  {\normalsize Word count: 8,753 \par}
  \vspace{1em}
  {\normalsize Data files used: O*NET 30.3; Annual Population Survey 2022--2025; with code submitted to \texttt{sgpe@ed.ac.uk}. \par}
\else
  {\large Tommy Ranyard \par}
  \vspace{0.5em}
  {\large MSc Economics, University of Edinburgh \par}
  \vspace{0.5em}
  {\large \today \par}
  \vspace{2cm}
  \begin{minipage}{0.85\textwidth}
  \small
  This dissertation is undertaken in collaboration with the Scottish Fiscal Commission (SFC) through the SGPE Sponsored Dissertation Programme. I thank my academic supervisors, Mark Schaffer and Eoin McLaughlin, along with Silvia Palombi and Alexander Sibelle at the SFC, for their helpful comments, and the OBR for insight into their approach, prompt, and results. All errors are my own.
  \end{minipage}
\fi

\vfill
\end{titlepage}

% ---------------------------------------------------------------------
\begin{abstract}
\noindent

The block grant adjustment responds to the difference between the Scottish and rUK tax bases and not to the level of either, so any shock to the Scottish fiscal position must act through that difference. Artificial intelligence is of the like since it will not affect every occupation equally, and what is uneven across occupations is uneven across regions. Yet the available exposure measures are point estimates, reported without addressing how much of the score came from the model that produced it. This dissertation establishes which parts of such a measure institutions can rely on, and how to audit economic claims built on it. The solution for this paper's setting comes rests on a coincidence that differencing across regions removes the common component of any model-specific bias. Standardised or rank-based differentials remove any surviving affine bias that rescales the index but does not reorder occupations. Scotland sits two percentile points below the rUK in the occupational exposure distribution, 0.86 index points on a 0--100 scale. Re-scoring all tasks on two further models preserves the sign and the percentile magnitude but not the index-point figure. The gap implies a ten-year TFP differential of three basis points and a net-position sensitivity near \pounds\fiscCen~million a year, whose largest uncertainty is the choice of scoring model. Exposure is no material differentiator in next year's Scottish forecast, though its direction and persistence bear on a system reconciled \textit{ex post}. Where magnitude does matter, the choice of model can shift the figures institutions lean on even more than the data do.
\end{abstract}

\medskip
\noindent\textbf{Keywords:} artificial intelligence, occupational exposure,
large language model scoring, model dependence, shift-share decomposition, Scotland.\\
\textbf{JEL classification:} C43, H68, H77, J24, O33, R23.

\clearpage